% ===== Section 3: "Achy Breaky Brain" (love, monogamy, break-ups) (~14 min) =====
\partaccent{brainpink}
\topicopener{brokenheart}{Սեր ու բաժանում}{Կոտրված սիրտ, ցավող ուղեղ}{Ինչու է բաժանումն այդքան ավերիչ}{sad}

% Heartbreak: nothing physical happened, so why does it hurt so much?
\begin{frame}{Բազմոցին, վարագույրները փակ}
  \begin{columns}[c]
    \begin{column}{0.44\textwidth}
      \begin{bookquote}
        Երբևէ օրերով պառկե՞լ ես բազմոցին կծկված, վարագույրները փակ, հեռախոսը չպատասխանած, ու մտածել, թե ինչու է ամբողջ տիեզերքը որոշել քեզ տանջել։
      \end{bookquote}
    \end{column}
    \begin{column}{0.52\textwidth}
      \begin{itemize}
        \item Ներշնչում է մեծ արվեստ ու երաժշտություն, ինչպես նաև \textbf{սարսափելի պոեզիա}
        \item<2-> Բայց ֆիզիկապես \textbf{ոչինչ չի պատահել}. ոչ վնասվածք, ոչ վիրուս
        \item<3-> Ինչու՞ է ամիսներով ցավում։ Որովհետև ուրիշները \textbf{մեծ ազդեցություն} ունեն ուղեղի բարօրության վրա
      \end{itemize}
    \end{column}
  \end{columns}
\end{frame}

% Culture is obsessed with couples; monogamy is odd for a primate
\begin{frame}{Ամբողջ մշակույթը զույգերի մասին է}
  \begin{columns}[c]
    \begin{column}{0.46\textwidth}
      \chip{heart}{Վալենտինի օր} \chip{couple}{հարսանիքներ}\par\smallskip
      \chip{film}{ռոմանտիկ կատակերգություններ}\par\smallskip
      \chip{musical}{սիրո բալլադներ} \chip{sparkles}{զարդեր}\par\smallskip
      \chip{memo}{պոեզիայի կեսը} \chip{loveletter}{բացիկներ}\par\medskip
      \onslide<3->{%
      \begin{chatbox}
        Գրիր բաժանման մասին ամենատխուր երգը, որ գիտես։
      \end{chatbox}}
    \end{column}
    \begin{column}{0.5\textwidth}
      \begin{itemize}
        \item Մոնոգամիան (մեկ զուգընկեր) մյուս պրիմատների մոտ \textbf{ընդունված չէ}
        \item<2-> Մենք ավելի երկար ենք ապրում ու կարող էինք շատ զուգընկերներ ունենալ։ Ավելի շատ զուգընկեր՝ ավելի շատ գեներ
        \item<2-> Բայց մարդիկ մեծ մասամբ ընտրում են \textbf{մեկին}
      \end{itemize}
    \end{column}
  \end{columns}
\end{frame}

% Why monogamy: three theories
\begin{frame}{Ինչու՞ մոնոգամիա. երեք վարկած}
  \begin{columns}[T]
    \begin{column}{0.31\textwidth}
      \begin{card}[colback=brainpink!25!cardwhite]
        \centering\bigemoji[0.8cm]{family}\par\smallskip
        \raggedright{\blackfont Երկու ծնող}\par
        Երկուսով խնամված երեխան ավելի հավանական է, որ գոյատևի
      \end{card}
    \end{column}
    \begin{column}{0.31\textwidth}
      \onslide<2->{%
      \begin{card}[colback=butter!45!cream]
        \centering\bigemoji[0.8cm]{moneybag}\par\smallskip
        \raggedright{\blackfont Մշակույթ}\par
        Հարստությունը պետք էր պահել ընտանիքում։ Իսկ ժառանգություն չես թողնի, եթե չգիտես՝ ով է քո ընտանիքը
      \end{card}}
    \end{column}
    \begin{column}{0.31\textwidth}
      \onslide<3->{%
      \begin{card}[colback=lavender!40!cardwhite]
        \centering\bigemoji[0.8cm]{baby}\par\smallskip
        \raggedright{\blackfont Տատիկներ}\par
        Տատիկներն օգնում էին թոռներին։ Բայց ոչ մի տատիկ չի խնամի երեխայի \textbf{նախկինի} երեխաներին
      \end{card}}
    \end{column}
  \end{columns}
  \vspace{3mm}
  \onslide<4->{\centering\small Ինչ էլ որ լինի պատճառը, սիրահարվելիս ուղեղը շատ տարօրինակ բաներ է անում։\par}
\end{frame}

% Attraction: antlers and tails
\begin{frame}{Գրավչություն. եղջյուրներ ու փետուրներ}
  \begin{columns}[c]
    \begin{column}{0.42\textwidth}
      \centering
      \includegraphics[width=\linewidth]{peacock.jpg}\par
      {\scriptsize Սիրամարգի պոչը գեղեցիկ է, բայց բազմանալու համար պետք չէ}
    \end{column}
    \begin{column}{0.54\textwidth}
      \begin{itemize}
        \item \textbf{Երկրորդային սեռական հատկանիշներ}. լոսի եղջյուրներ, սիրամարգի պոչ։ Ցույց են տալիս, թե որքան առողջ է կենդանին
        \item<2-> Մարդիկ էլ ունեն. խորը ձայն ու մորուք, կուրծք ու կորություններ
        \item<3-> Ինչն էր առաջինը՝ հատկանիշը, թե՞ ուղեղի սերը դրա նկատմամբ։ \textbf{Հավի ու ձվի} հարց
      \end{itemize}
    \end{column}
  \end{columns}
\end{frame}

% Preferences: wit, culture, and people who look like us
\begin{frame}{Ճաշակի հարց}
  \begin{columns}[T]
    \begin{column}{0.31\textwidth}
      \begin{card}[colback=mint!35!cardwhite]
        \centering\bigemoji[0.9cm]{bulb}\par\smallskip
        \raggedright{\blackfont Խելք}\par
        Ոմանց համար ամենագրավիչը մարդու սրամտությունն է
      \end{card}
    \end{column}
    \begin{column}{0.31\textwidth}
      \onslide<2->{%
      \begin{card}[colback=butter!45!cream]
        \centering\bigemoji[0.9cm]{earth}\par\smallskip
        \raggedright{\blackfont Մշակույթ}\par
        Արևմուտքում՝ արհեստական արևայրուք, Ասիայում՝ մաշկը \textbf{սպիտակեցնող} քսուքներ
      \end{card}}
    \end{column}
    \begin{column}{0.31\textwidth}
      \onslide<3->{%
      \begin{card}[colback=brainpink!25!cardwhite]
        \centering\bigemoji[0.9cm]{mirror}\par\smallskip
        \raggedright{\blackfont Նմանություն}\par
        Մեզ ավելի են գրավում \textbf{մեզ նման} մարդիկ։ Ուղեղի եսասիրությունը նորից
      \end{card}}
    \end{column}
  \end{columns}
\end{frame}

% Lust is not love: Bartels & Zeki
\begin{frame}{Ցանկություն $\neq$ սեր}
  \relax{\small Սեռական ցանկությունն (lust) ու սերը կոնկրետ մեկի հանդեպ ուղեղը մշակում է \textbf{տարբեր} կերպ։ Բարտելսն ու Զեկին սիրահարվածներին ցույց են տվել զուգընկերոջ նկարը.}\par\medskip
  \begin{columns}[T]
    \begin{column}{0.34\textwidth}
      \begin{card}[colback=mint!35!cardwhite]
        {\blackfont $\uparrow$ Ակտիվանում են}\par\smallskip
        կղզյակ (insula)\\
        գոտևոր կեղևի առաջային մաս\\
        պոչավոր կորիզ, պուտամեն
      \end{card}
    \end{column}
    \begin{column}{0.34\textwidth}
      \onslide<2->{%
      \begin{card}[colback=lavender!40!cardwhite]
        {\blackfont $\downarrow$ Հանգստանում են}\par\smallskip
        գոտևոր կեղևի հետին մաս՝ \textbf{ցավ}\\
        նշաձև մարմին՝ \textbf{վախ, զայրույթ}\\
        նախաճակատային կեղև՝ \textbf{տրամաբանություն}
      \end{card}}
    \end{column}
    \begin{column}{0.26\textwidth}
      \onslide<3->{%
      \centering\mascot[1.5cm]{smug}\par\smallskip
      \raggedright\small Դրա համար են զույգերն այդքան հանգիստ ու \textbf{ինքնագոհ} երևում}
    \end{column}
  \end{columns}
\end{frame}

% Dopamine: love works like a drug (meme)
\begin{frame}{Սերը՝ ինչպես թմրանյութ}
  \begin{columns}[c]
    \begin{column}{0.42\textwidth}
      \begin{itemize}
        \item Սիրահարվածությունը բարձրացնում է \textbf{դոֆամինի} ակտիվությունը պարգևի համակարգում
        \item Զուգընկերոջ ներկայությունը հաճույք է, \textbf{գրեթե թմրանյութի պես}
        \item Կախվածության մասին՝ 8-րդ գլխում
      \end{itemize}
    \end{column}
    \begin{column}{0.54\textwidth}
      \centering
      \onslide<2->{%
      \begin{meme}[7cm]{epic_handshake.jpg}
        \memecap{0.17}{0.4}{Սիրահարվածություն}
        \memecap{0.8}{0.55}{Թմրանյութեր}
        \memecap[font=\blackfont\large]{0.4}{0.9}{Դոֆամին}
      \end{meme}}
    \end{column}
  \end{columns}
\end{frame}

% Oxytocin and pheromones
\begin{frame}{Օքսիտոցին ու ֆերոմոններ}
  \begin{columns}[T]
    \begin{column}{0.5\textwidth}
      \begin{itemize}
        \item \textbf{Օքսիտոցին}՝ «սիրո հորմոն»։ Բըրնեթ. բարդ նյութի \textbf{ծիծաղելի պարզեցում} է։ Բարձր է զույգերի մոտ ու կապվում է վստահության հետ
      \end{itemize}
      \onslide<2->{%
      \begin{scienceupdate}
        Վստահության մասին պնդումը լավ չի կրկնվում. քթի սփրեյով օքսիտոցինի ազդեցությունը հետագա փորձերում չհաստատվեց։
      \end{scienceupdate}}
    \end{column}
    \begin{column}{0.46\textwidth}
      \onslide<3->{%
      \begin{card}[colback=brainpink!25!cardwhite]
        {\blackfont Ֆերոմոններ}\par\smallskip
«Ֆերոմոններով» սփրեյներ կան, բայց մարդկային «գրավչության» ֆերոմոնների \textbf{ապացույց չկա}։
      \end{card}}
      \onslide<4->{%
      \begin{bookquote}
        Ուղեղը հաճախ ապուշ է, բայց այդքան հեշտ էլ չի կառավարվում։
      \end{bookquote}}
    \end{column}
  \end{columns}
  \source{Nave, Camerer \& McCullough, Perspect. Psychol. Sci., 2015}
\end{frame}

% The partner becomes part of you
\begin{frame}{Զուգընկերը դառնում է քո մի մասը}
  \begin{columns}[c]
    \begin{column}{0.58\textwidth}
      \begin{itemize}
        \item Մի ամբողջ մարդ քեզ այդքան գնահատում է. մեծ \textbf{ձեռքբերում} է
        \item Զույգ լինելը սոցիալական կարգավիճակ է տալիս (ինչպես կասի ցանկացած երջանիկ միայնակ մարդ՝ \textbf{ատամները սեղմած})
        \item<2-> Ուղեղը հարմարվում է ու սկսում է այս ամենը \textbf{սպասել}. զուգընկերը մտնում է քո ծրագրերի ու նպատակների մեջ
      \end{itemize}
    \end{column}
    \begin{column}{0.38\textwidth}
      \centering
      \begin{tikzpicture}
        \fill[brainpink!35!cardwhite] (0,0) circle (1.2);
        \fill[lavender!55!cardwhite, visible on=<1>] (2.1,0) circle (1.2);
        \fill[lavender!55!cardwhite, visible on=<2->] (1.3,0) circle (1.2);
        \begin{scope}[visible on=<2->]
          \clip (0,0) circle (1.2);
          \fill[blush!70!cardwhite] (1.3,0) circle (1.2);
        \end{scope}
        \draw[plum, line width=0.9pt] (0,0) circle (1.2);
        \draw[plum, line width=0.9pt, visible on=<1>] (2.1,0) circle (1.2);
        \draw[plum, line width=0.9pt, visible on=<2->] (1.3,0) circle (1.2);
        \node[font=\blackfont, text=plum] at (-0.55,0) {Ես};
        \node[font=\blackfont, text=plum, visible on=<1>] at (2.6,0) {Նա};
        \node[font=\blackfont, text=plum, visible on=<2->] at (1.9,0) {Նա};
        \node[font=\blackfont\small, text=plum, visible on=<2->] at (0.65,0) {Մենք};
      \end{tikzpicture}
    \end{column}
  \end{columns}
\end{frame}

% And then it ends
\begin{frame}{Եվ հետո այն ավարտվում է}
  \begin{columns}[T]
    \begin{column}{0.46\textwidth}
      \begin{itemize}
        \item Անհավատարմություն, անհամատեղելիություն, մեկի վարքը վանեց մյուսին
        \item Տագնապային մարդիկ հակված են \textbf{ուռճացնել} կոնֆլիկտները, երբեմն՝ մինչև խզում
        \item<2-> Ինչքան բան էր ուղեղը ներդրել։ Ու այդ ամենը վերանում է \textbf{մեկ հարվածով}
      \end{itemize}
    \end{column}
    \begin{column}{0.5\textwidth}
      \onslide<3->{%
      \begin{card}[colback=brainpink!25!cardwhite]
        {\blackfont Ուղեղը միանգամից կորցնում է}\par\smallskip
        \emoji{gift}\ սպասվող բոլոր հաճելի զգացումները\par\smallskip
        \emoji{calendar}\ ապագայի ծրագրերը։ Իսկ ուղեղը \textbf{անորոշություն} չի սիրում\par\smallskip
        \emoji{house}\ գործնական պատասխանները. որտե՞ղ ապրեմ, ընկերներին կկորցնե՞մ
      \end{card}}
    \end{column}
  \end{columns}
\end{frame}

% Social pain, literally
\begin{frame}{Սա իսկապես ցավում է}
  \begin{columns}[T]
    \begin{column}{0.52\textwidth}
      \begin{itemize}
        \item Մեկը, ով քեզ ճանաչում է բոլորից լավ, որոշել է, որ \textbf{պետք չես}։ Հարված սոցիալական ինքնությանը
        \item<2-> Բաժանումը ակտիվացնում է այն շրջանները, որոնք մշակում են \textbf{ֆիզիկական ցավը}։ «Սերը ցավ է պատճառում»՝ բառացիորեն
        \item<2-> Պարացետամոլը երբեմն նույնիսկ օգնում է «սրտի ցավին»
      \end{itemize}
    \end{column}
    \begin{column}{0.44\textwidth}
      \onslide<3->{%
      \begin{scienceupdate}
        2014-ին նույն տվյալները վերլուծեցին ավելի նուրբ մեթոդով։ Շրջանները մասամբ համընկնում են, բայց ակտիվության \textbf{նախշերը տարբեր} են։ Նման է ցավին, բայց \textbf{նույնը չէ}։
      \end{scienceupdate}}
    \end{column}
  \end{columns}
  \source{Kross et al., PNAS, 2011; Woo et al., Nat. Commun., 2014}
\end{frame}

% Memories turn sour; withdrawal
\begin{frame}{Հիշողություններ ու «զրկանք»}
  \begin{columns}[c]
    \begin{column}{0.58\textwidth}
      \begin{itemize}
        \item Անհամար երջանիկ հիշողություններ հիմա կապված են \textbf{ցավի} հետ։ Սա խարխլում է քո «ես»-ի մի մեծ մասը
        \item<2-> Սովոր էիր մշտական պարգևի, ու այն հանկարծ վերցրին
        \item<3-> Մի բան, որ նման է թմրանյութից \textbf{զրկվելուն} (withdrawal)։ Ավելին՝ 8-րդ գլխում
      \end{itemize}
    \end{column}
    \begin{column}{0.38\textwidth}
      \centering
      \bigemoji[1cm]{camera}\ \bigemoji[1cm]{loveletter}\ \bigemoji[1cm]{musical}\par\medskip
      {\Large$\downarrow$}\par\medskip
      \bigemoji[1.4cm]{brokenheart}
    \end{column}
  \end{columns}
\end{frame}

% Recovery: it does get better
\begin{frame}{Ժամանակը իսկապես բուժում է}
  \begin{columns}[c]
    \begin{column}{0.58\textwidth}
      \begin{itemize}
        \item Ուղեղը կարողանում է ամեն ինչ \textbf{նորից հավաքել}, թեև դանդաղ
        \item<2-> Երբ կենտրոնանում ես բաժանման \textbf{դրական հետևանքների} վրա (օրինակ՝ գրելով), ավելի արագ ես վերականգնվում
        \item<3-> Գիտությունն ու կլիշեն համընկնում են. ժամանակի հետ \textbf{ավելի լավ է դառնում}
      \end{itemize}
      \onslide<4->{%
      \begin{bookquote}
        «Բաժանվելը դժվար է» ասելը շատ մեղմ է ասված։
      \end{bookquote}}
    \end{column}
    \begin{column}{0.38\textwidth}
      \centering
      \bigemoji[1.2cm]{sunrise}\par\medskip
      \mascot[2.3cm]{happy}
    \end{column}
  \end{columns}
\end{frame}
